\documentclass[11pt]{article}

\usepackage[T1]{fontenc}
\usepackage[margin=1in]{geometry}
\usepackage{amsmath,amssymb,amsthm,mathtools}
\usepackage{hyperref}
\usepackage{enumitem}

\newtheorem{theorem}{Theorem}
\newtheorem{lemma}[theorem]{Lemma}
\newtheorem{proposition}[theorem]{Proposition}
\newtheorem{remark}[theorem]{Remark}

\newcommand{\Pset}{\mathbb P}
\newcommand{\R}{\mathcal R}
\newcommand{\B}{\mathcal B}
\newcommand{\C}{\mathcal C}
\newcommand{\Hcal}{\mathcal H}
\newcommand{\eps}{\varepsilon}

\title{A proof of Erd\H{o}s Problem 689}
\author{Working note}
\date{April 26, 2026 (revised)}

\begin{document}
\maketitle

\begin{abstract}
We prove Erd\H{o}s Problem 689 for all sufficiently large \(n\): for every
sufficiently large \(n\) there is a choice of congruence classes
\(a_p\pmod p\) for all primes \(p\le n\) such that every integer in \([1,n]\)
satisfies at least two of the congruences \(\equiv a_p\pmod p\).  The proof
uses a fixed small-prime sieve, robust cleanup primes larger than \(n/5\), a
finite-core prime-difference matching, the Green--Tao finite-complexity
linear-forms theorem in primes, made unconditional using the Green--Tao
nilsequence theorem and the Green--Tao--Ziegler inverse theorem for Gowers
norms, and Kahn's fractional Frankl--R\"odl--Pippenger rounding theorem.

The application of Kahn's Theorem 1.5 has been verified directly against the
printed paper (\emph{Random Structures and Algorithms} 8 (1996), 149--157):
the $C(e)\equiv 1$ application falls inside the printed hypothesis as a
3-bounded codegree-2 fractional matching with $\sum_e t_e\to\infty$, giving
the conclusion $|M_n|=(1-o(1))|Z_n|$.  The application of Green--Tao--Ziegler
is reduced to the standing finite-complexity theorem applied to four explicit
affine-linear systems whose coefficient vectors are written out in
Appendix~B; admissibility and diagonal deletion are verified in
Appendices~C--D.

This note was prepared with substantial AI assistance and should be
independently checked.
\end{abstract}

\paragraph{Relation to the existing EP689 discussion.}
The discussion thread for EP689 already points toward this technology: earlier
comments suggest using linear equations in primes to control averaged
prime-pattern counts and Pippenger--Spencer/Kahn-style hypergraph covering to
turn fractional coverage into a genuine disjoint cover.  The new feature of
the present note is the robust \(P>n/5\) cleanup setup, the finite-core
prime-difference matching formulation with explicit half-residue kernel, and
the strict surplus margin
\(\Delta=(\beta+\tfrac35)\delta_S-1>0\) which absorbs all coefficient tails
and exceptional residual tokens.

\section{Statement and conventions}

For a prime \(p\le n\), write \(a_p \pmod p\) for the residue class chosen for
\(p\).  The goal is to choose the \(a_p\)'s so that every integer
\(m\in[1,n]\) is hit by at least two chosen congruence classes:
\[
  \#\{p\le n:\ m\equiv a_p \pmod p\}\ge 2.
\]
We leave most primes in their zero class.  The prime \(2\) is put in the odd
class \(a_2\equiv1\pmod 2\), and \(3\) is left at \(0\pmod 3\).  A fixed finite
set \(S\subset\{7,11,13,\ldots\}\) is switched to nonzero residues
\(b_s\pmod s\).  Further primes \(P>n/5\) will be used as cleanup primes.

For the fixed set \(S\), put
\[
  H_S(x)=\#\{s\in S:\ x\equiv b_s\pmod s\}.
\]
After the fixed switches in \(S\), the main residual demand consists of even
integers of the form
\[
  x=2^k u q,\qquad k\ge1,
\]
where \(u\) is odd \(S\)-smooth and \(q\notin S\) is an odd prime appearing to
exponent one, subject to the finite system of congruence exclusions
\(H_S(x)=0\).  Prime powers, pure
\(\{2\}\cup S\)-smooth terms, and small-prime boundary cases contribute only
\(o(n/\log n)\) residual tokens.  Standard fixed-modulus prime number theory
and the identity
\[
  \sum_{k\ge1}2^{-k}
  \sum_{A\subseteq S}
  \prod_{s\in A}\frac1{s-1}
  \prod_{s\notin A}\left(1-\frac1{s-1}\right)=1
\]
give
\[
  |A_S(n)|=(1+o(1))\frac n{\log n}.
\]
Moreover, if
\[
  A_1(n)=\{x\in A_S(n):v_2(x)=1\},\qquad
  A_2(n)=\{x\in A_S(n):v_2(x)\ge2\},
\]
then
\[
  |A_1(n)|=\left(\frac12+o(1)\right)|A_S(n)|,\qquad
  |A_2(n)|=\left(\frac12+o(1)\right)|A_S(n)|.
\]

\section{Robust cleanup primes}

Let
\[
  W=\prod_{s\in S}s.
\]
For a unit residue \(\pi\pmod W\), call \(\pi\) robust if
\[
  H_S(\pi)\ge1,\qquad H_S(2\pi)\ge2,\qquad H_S(4\pi)\ge2.
\]
Let \(\B\subset(\mathbb Z/W\mathbb Z)^\times\) be the set of robust classes and
\[
  \delta_S=\frac{|\B|}{\varphi(W)}.
\]
For each \(s\in S\), the three classes
\[
  \pi\equiv b_s,\qquad 2\pi\equiv b_s,\qquad 4\pi\equiv b_s\pmod s
\]
are distinct.  Thus \(\delta_S\) depends only on \(S\), not on the particular
nonzero residues \(b_s\).  A union bound gives
\[
  \delta_S\ge 1-A_S^{(0)}(3+2B_S),
\]
where
\[
  A_S^{(0)}=\prod_{s\in S}\left(1-\frac1{s-1}\right),\qquad
  B_S=\sum_{s\in S}\frac1{s-2}.
\]
For instance, take \(S(y)=\{p:\ 7\le p\le y,\ p\text{ prime}\}\).  By Mertens'
theorem and the prime number theorem,
\[
  A_{S(y)}^{(0)}\asymp\frac1{\log y},\qquad
  B_{S(y)}=O(\log\log y),
\]
so \(A_{S(y)}^{(0)}(3+2B_{S(y)})\to0\).  Hence we may choose a fixed finite
\(S\) so that \(\delta_S\) is as close to \(1\) as needed.  This \(S\) may be
very large; its size is fixed before the limit \(n\to\infty\).

Set
\[
  \beta_*=\frac12\left(1-\frac35e^{-2}\right)=0.459399\ldots
\]
and
\[
  \delta_*=\frac1{\beta_*+3/5}=0.94393\ldots .
\]
Choose \(S\) so that \(\delta_S>\delta_*\), and then choose
\[
  \delta_S^{-1}-\frac35<\beta<\beta_*.
  \tag{1}
\]
Let
\[
  \Delta=(\beta+\tfrac35)\delta_S-1>0.
  \tag{2}
\]
This fixed positive margin is the surplus that later absorbs coefficient tails
and the \(o(n/\log n)\) exceptional residual tokens.
Let
\[
  \R_\beta(n)=\{P\in(n/5,\beta n]: P\text{ prime},\ P\bmod W\in\B\}.
\]
By the prime number theorem in fixed arithmetic progressions,
\[
  |\R_\beta(n)|=\left((\beta-\tfrac15)\delta_S+o(1)\right)\frac n{\log n}.
\]
The full robust reservoir
\[
  \R(n)=\{P\in(n/5,n]:P\text{ robust}\}
\]
has size
\[
  |\R(n)|=\left(\frac45\delta_S+o(1)\right)\frac n{\log n}.
\]

\begin{lemma}[Robust side-debt]
Switching a robust prime \(P>n/5\) to any nonzero residue creates no unresolved
side debt.
\end{lemma}

\begin{proof}
The only multiples of \(P\le n\) are \(P,2P,3P,4P\).  The integer \(P\) is
covered by the parity class and at least one \(S\)-hit, since \(H_S(P)\ge1\).
The integers \(2P\) and \(4P\) have at least two \(S\)-hits by robustness.  The
integer \(3P\) is covered by the parity class and the unchanged zero class
modulo \(3\).  There is no multiple \(5P\le n\).
\end{proof}

Consequently, switching any collection of robust primes \(P>n/5\) to nonzero
residues creates no new residual tokens: the only zero-class hit lost at a
multiple of such a \(P\) is replaced by two hits not using the zero class of
that \(P\).

\section{The matching hypergraph}

Choose finite coefficient cores
\[
  \mathcal A_1\subset\{a:a\text{ odd }S\text{-smooth}\},
\]
\[
  \mathcal A_2\subset\{b:b=2^j u,\ j\ge1,\ u\text{ odd }S\text{-smooth}\}.
\]
The core \(X_n\subset A_1(n)\) consists of residual targets \(x=2aq\) with
\(a\in\mathcal A_1\).  The core \(Y_n\subset A_2(n)\) consists of residual
targets \(y=2bq'\) with \(b\in\mathcal A_2\).  The finite cores are chosen to
capture fractions \(\alpha_X,\alpha_Y\) of the \(A_1,A_2\) coefficient mass,
uniformly in the half-residue classes defined below, with
\[
  \alpha_X,\alpha_Y>G(\beta)+\eta
\]
for some fixed \(\eta>0\), where
\[
  G(\beta)=\int_{1/5}^{\beta}\frac{dt}{1-2t}
  =\frac12\log\left(\frac{3/5}{1-2\beta}\right)<1.
  \tag{3}
\]
This is possible because the full \(S\)-smooth coefficient distributions are
absolutely summable.  The order of quantifiers is: after fixing \(S\), \(\beta\),
and \(\Delta\), choose the coefficient cores so that the omitted main residual
coefficient share is smaller than a fixed fraction of \(\Delta\) and so that the
above side-load inequalities hold.  Then take \(n\) large enough that the actual
coefficient-tail counts, the exceptional residual-token count, and all GTZ/Kahn
losses are smaller than the remaining margin.

Let \(Z_n=\R_\beta(n)\).  Define a 3-partite 3-uniform hypergraph
\[
  \Hcal_n\subset X_n\times Y_n\times Z_n
\]
by putting an edge \((x,y,P)\) whenever
\[
  |y-x|=2P.
\]
The parity split is essential: if both targets had the same 2-adic parity
class, then \((y-x)/2\) would be even and could not equal the odd prime \(P\).

The goal is to find a matching in \(\Hcal_n\) covering
\[
  (1-o(1))|Z_n|
\]
labels.  Such a matching covers that many pairs of residual targets using one
robust prime per pair.

\section{The half-residue transport}

Put
\[
  c_s\equiv 2^{-1}b_s\pmod s
\]
and define the set of residual half-residues
\[
  \C=\{A\bmod W:\ A\not\equiv c_s\pmod s\text{ for every }s\in S\}.
\]
For \(x=2aq\), write \(A\equiv aq\pmod W\); for \(y=2bq'\), write
\(B\equiv bq'\pmod W\).  Then residual membership is simply
\[
  A,B\in\C.
\]
For every unit \(\pi\pmod W\) and every orientation \(\sigma=\pm1\),
\[
  \#\{A\in\C:\ A+\sigma\pi\in\C\}=M:=\prod_{s\in S}(s-2).
  \tag{4}
\]
Locally this just excludes two distinct residues modulo each \(s\in S\), and
they are distinct because \(\pi\) is a unit.

In the aggregate continuum model, a label is \((t,\pi)\) with
\[
  t=P/n\in(1/5,\beta],\qquad \pi\in\B.
\]
For each label, choose \(\sigma=\pm1\) with equal probability, choose uniformly
an \(A\in\C\) with \(A+\sigma\pi\in\C\), and choose \(u\in(0,1-2t)\) uniformly.
For \(\sigma=+1\), send mass to
\[
  X=(u,A),\qquad Y=(u+2t,A+\pi),
\]
and for \(\sigma=-1\), send mass to
\[
  X=(u+2t,A),\qquad Y=(u,A-\pi).
\]
This gives exact label load \(L_Z(t,\pi)=1\).  The side load is bounded by
\[
  L_X,L_Y\le \frac{G(\beta)}{\alpha_X}+o(1),\qquad
  \frac{G(\beta)}{\alpha_Y}+o(1),
\]
so by the choice of core it is at most \(1-2\gamma\) for some fixed
\(\gamma>0\).

\section{Typed-kernel lift}

A type is a tuple
\[
  \tau=(a,b,\sigma,r,r',\pi),
\]
where \(a\in\mathcal A_1\), \(b\in\mathcal A_2\), \(\sigma\in\{\pm1\}\),
\(r,r'\in(\mathbb Z/W\mathbb Z)^\times\), and \(\pi\in\B\) satisfy
\(B:=br'\equiv ar+\sigma\pi=A+\sigma\pi\pmod W\), with \(A,B\in\C\).
Its support polygon is
\[
  D_\tau=\left\{(Q,Q'):0<Q\le\frac1{2a},\
  0<Q'\le\frac1{2b},\
  \frac15<\sigma(bQ'-aQ)\le\beta\right\}.
\]

\subsection*{Half-residue marginal identity}

For \(A\in\C\), define
\[
  \Xi_A^\infty
  =\sum_{\substack{a\text{ odd }S\text{-smooth}\\ ar\equiv A\,(W)}}\frac1{2a},
  \qquad
  H_A^\infty
  =\sum_{\substack{b=2^ju,\,j\ge1,\,u\text{ odd }S\text{-smooth}\\
                   br'\equiv A\,(W)}}\frac1{2b}.
\]

\begin{lemma}[Uniform half-residue marginals]
For every \(A\in\C\),
\[
  \Xi_A^\infty=H_A^\infty=\tfrac12.
\]
\end{lemma}

\begin{proof}
Work locally at \(s\in S\).  For the \(X\)-side, if \(A\not\equiv0\pmod s\),
then \(v_s(a)=0\) and there is exactly one unit residue \(r_s\) with
\(ar_s\equiv A\pmod s\), giving local mass~\(1\).  If \(A\equiv0\pmod s\),
then \(v_s(a)=e\ge1\) and any of the \(s-1\) unit residues is possible,
giving \(\sum_{e\ge1}(s-1)s^{-e}=1\).  The product over \(s\in S\) is~\(1\),
and the global \(\tfrac12\) in \(1/(2a)\) supplies the displayed value.  The
\(Y\)-side has the same \(S\)-local computation, and its two-adic factor is
\(\sum_{j\ge1}2^{-j-1}=\tfrac12\).
\end{proof}

The key consequence is that \(\Xi_A^\infty\) and \(H_A^\infty\) are
\emph{independent of} \(A\in\C\): they are equal to the same constant
\(\tfrac12\) for every residual half-residue.  Finite cores
\(\mathcal A_1,\mathcal A_2\) capture fractions
\(\alpha_X(A)=\Xi_A/\Xi_A^\infty\) and \(\alpha_Y(A)=H_A/H_A^\infty\) of these
limits; setting \(\alpha_X=\min_A\alpha_X(A)\) and
\(\alpha_Y=\min_A\alpha_Y(A)\), absolute summability lets us pick cores with
\(\alpha_X,\alpha_Y\) arbitrarily close to \(1\), hence in particular larger
than \(G(\beta)\).

\subsection*{Construction of typed kernels}

For \(A\in\C\), define the conditional disintegrations
\[
  \rho_X^A(a,r)=\frac{1/(2a)}{\Xi_A},\qquad
  \rho_Y^B(b,r')=\frac{1/(2b)}{H_B},
\]
which are probability measures on the relevant finite-core fibers.

For \((Q,Q')\in D_\tau\), let
\[
  t=\sigma(bQ'-aQ)\in(1/5,\beta],
  \qquad
  u_\tau(Q,Q')=
  \begin{cases}2aQ,&\sigma=+1,\\ 2bQ',&\sigma=-1.\end{cases}
\]
The map \((Q,Q')\mapsto(u_\tau,t)\) on \(D_\tau\) has absolute Jacobian
\(2ab\).  Define the aggregate kernel
\[
  K_\sigma(t,\pi;u,A)=\frac1{2M(1-2t)}
  \quad\text{on }0<u<1-2t,\,A\in\C,\,A+\sigma\pi\in\C,
\]
which carries unit mass per label by~(4).  The pre-normalized typed intensity
on \(D_\tau\) is
\[
  h_\tau(Q,Q')
  =2ab\,K_\sigma(t,\pi;u_\tau,A)\,\rho_X^A(a,r)\rho_Y^B(b,r').
\]
Let \(\kappa_\tau>0\) be the fixed local GTZ constant normalized so that, for
bounded piecewise-continuous \(F\) supported in \(D_\tau\),
\[
  \sum_{\substack{q\equiv r,\,q'\equiv r'\,(W)\\ q,q',\sigma(bq'-aq)\text{ prime}\\
                   \sigma(bq'-aq)\equiv\pi\,(W)}}
  \frac{\log^2 n}{n}F(q/n,q'/n)
  =
  \left(\frac{n}{\varphi(W)\log n}+o\!\left(\frac n{\log n}\right)\right)
  \kappa_\tau\int_{D_\tau}F\,dQ\,dQ'.
\]
Set \(g_\tau=h_\tau/\kappa_\tau\).  Since the type set is finite,
\(t\le\beta<1/2\), and \(\alpha_X,\alpha_Y>0\), the kernels \(g_\tau\) are
uniformly bounded on \(D_\tau\).

\begin{proposition}[Typed-kernel lift]
The typed kernels \(g_\tau\) satisfy
\[
  L_Z^{\rm lim}(\pi,t)=1
  \quad\text{for every }(\pi,t)\in\B\times(1/5,\beta],
\]
\[
  L_X^{\rm lim}(a,r,Q)\le\frac{G(\beta)}{\alpha_X},\qquad
  L_Y^{\rm lim}(b,r',Q')\le\frac{G(\beta)}{\alpha_Y}.
\]
In particular, after choosing cores with \(\alpha_X,\alpha_Y>G(\beta)\),
\[
  L_X^{\rm lim},L_Y^{\rm lim}\le 1-2\gamma
\]
for some fixed \(\gamma>0\).
\end{proposition}

\begin{proof}
\emph{Label load.}  Fix \((\pi,t)\) and let \(\sigma=+1\) (the
\(\sigma=-1\) case is identical).  On the line \(t=bQ'-aQ\),
\(Q'=(aQ+t)/b\), with coarea factor \(dQ/b\).  Then with \(u=2aQ\),
\[
  \frac1bh_\tau(Q,(aQ+t)/b)\,dQ
  =K_+(t,\pi;u,A)\,\rho_X^A(a,r)\rho_Y^{A+\pi}(b,r')\,du.
\]
Summing over \((a,r)\) above a fixed \(A\) gives \(\sum_{(a,r):ar=A}\rho_X^A=1\),
and summing over \((b,r')\) above \(B=A+\pi\) gives \(1\).  Hence
\[
  L_Z^{\rm lim}(\pi,t)
  =\sum_{\sigma}\sum_{A\in\C,A+\sigma\pi\in\C}\int K_\sigma\,du
  =1
\]
by~(4).

\emph{Side load.}  Fix \((a,r,Q)\) and put \(z=2aQ\), \(A=ar\).  Summing the
typed densities incident to this \(X\)-vertex gives
\[
  L_X^{\rm lim}(a,r,Q)=\frac{\Lambda_X^{\rm raw}(z,A)}{\Xi_A},
\]
where
\(
  \Lambda_X^{\rm raw}(z,A)
  =\frac{N_+(A)}{2M}J_\beta((1-z)/2)+\frac{N_-(A)}{2M}J_\beta(z/2)
\)
with \(J_\beta(s)=\int_{1/5}^{\min(\beta,s)}dt/(1-2t)\) and
\(N_\pm(A)\le M\).  Thus
\[
  L_X^{\rm lim}(a,r,Q)
  \le\frac{1}{\alpha_X(A)}
  \left[J_\beta\!\left(\frac{1-z}{2}\right)+J_\beta\!\left(\frac{z}{2}\right)\right].
\]
For \(\beta\ge3/10\), the bracketed sum is at most \(G(\beta)\): if
\(\max((1-z)/2,z/2)\ge\beta\) then the other term is at most
\(J_\beta(1/5)=0\); if both lie in \((1/5,\beta)\), they lie in
\((1/5,3/10)\), and convexity of \((1-2t)^{-1}\) gives the bound at the
endpoint.  Hence \(L_X^{\rm lim}\le G(\beta)/\alpha_X\).  The \(Y\)-side is
symmetric.
\end{proof}

\section{GTZ averaged moments}

For an edge \(e=(x,y,P)\) of type \(\tau\), with
\[
  x=2aq,\qquad y=2bq',\qquad P=\sigma(bq'-aq),
\]
assign weight
\[
  w_e=\frac{\log^2 n}{n}g_\tau(q/n,q'/n).
\]
We use the Green--Tao finite-complexity theorem for affine-linear forms in
primes, made unconditional using the Green--Tao nilsequence theorem and the
Green--Tao--Ziegler inverse theorem for Gowers norms.  The fixed modulus
\(2W\) is part of the residue data.  In Appendix~A we record the normalized
\(W_{\rm GTZ}\)-tricked formulation in which all subsequent moments are
stated; the four explicit affine-linear systems are written out in
Appendix~B; diagonal and exceptional-set deletions are handled in
Appendix~C; local admissibility is checked in Appendix~D.  The coefficient
cores, residue data, and kernels are fixed before \(n\to\infty\).

\begin{proposition}[Weighted moment proposition]
The loads
\[
  L_Z(P)=\sum_{e\ni P}w_e,\qquad
  L_X(x)=\sum_{e\ni x}w_e,\qquad
  L_Y(y)=\sum_{e\ni y}w_e
\]
satisfy
\[
  \sum_{P\in Z_n}L_Z(P)=|Z_n|+o(|Z_n|),
\]
\[
  \sum_{P\in Z_n}(L_Z(P)-1)^2=o(|Z_n|),
\]
\[
  \sum_{x\in X_n}(L_X(x)-L_X^{\rm lim}(x))^2=o(|X_n|),
\]
and
\[
  \sum_{y\in Y_n}(L_Y(y)-L_Y^{\rm lim}(y))^2=o(|Y_n|),
\]
where the limiting side loads are bounded by \(1-2\gamma\).  Moreover
\[
  \max_e w_e\ll \frac{\log^2 n}{n}=o(1).
\]
\end{proposition}

\begin{proof}
This is the conclusion of Proposition~F.1, proved via the systems of
Appendix~B, the deletions of Appendix~C, and the local admissibility checks
of Appendix~D.  No pointwise estimate for fixed \(P=bq'-aq\) is used; every
second moment retains the shared vertex as an averaging variable
(Appendix~E).
\end{proof}

\section{Fractional rounding}

Normalize the label loads by replacing \(w_e\) with \(w'_e=c_Pw_e\), where
\[
  c_P=
  \begin{cases}
    \min(1,L_Z(P)^{-1}),& L_Z(P)>0,\\
    1,& L_Z(P)=0
  \end{cases}
\]
for the unique label \(P\in e\).  The \(L^2\) estimate on \(Z_n\) gives
\[
  \sum_e w'_e=\sum_{P\in Z_n}\min(L_Z(P),1)=|Z_n|-o(|Z_n|).
\]
Delete all edges incident to side vertices with normalized side load \(>1\).
Let
\[
  B_X=\{x:\ L'_X(x)>1\}.
\]
Since \(L'_X(x)\le L_X(x)\) and \(L_X^{\rm lim}(x)\le1-2\gamma\), every
\(x\in B_X\) satisfies
\(|L_X(x)-L_X^{\rm lim}(x)|\ge 2\gamma\).  Thus the \(L^2\) estimate gives
\(|B_X|=o(|X_n|)\).  Moreover
\[
  \sum_{x\in B_X}L'_X(x)\le \sum_{x\in B_X}L_X(x)
  \le (1-2\gamma)|B_X|
  + |B_X|^{1/2}
    \left(\sum_x (L_X(x)-L_X^{\rm lim}(x))^2\right)^{1/2}
  =o(|Z_n|).
\]
The same argument applies to the \(Y\)-side.  Hence deleting overloaded side
vertices loses only \(o(|Z_n|)\) total mass.  The remaining weights
\(t_e\le w_e\) form a fractional matching:
\[
  \sum_{e\ni v}t_e\le1\qquad\text{for every vertex }v,
\]
with
\[
  \sum_e t_e=(1-o(1))|Z_n|,\qquad
  \max_e t_e\le\max_e w_e\ll \frac{\log^2 n}{n}=o(1).
\]
The hypergraph has codegree at most \(2\): a pair \((x,y)\) determines
\(P=|y-x|/2\), and a pair \((x,P)\) or \((y,P)\) has at most two extensions.
Therefore the fractional pair co-load satisfies
\[
  \alpha(t)=\max_{u\ne v}\sum_{e\supset\{u,v\}}t_e
  \le2\max_e t_e=o(1).
\]

By Kahn's fractional Frankl--R\"odl--Pippenger theorem (Theorem~1.5 of
\cite{Kahn1996}), applied with the single statistic \(C(e)\equiv1\), this
fractional matching rounds to an actual matching \(M_n\).  We verify the
hypothesis: with \(C\equiv 1\), the printed condition
\[
  \sum_e C(e)^2 t_e<o\!\left(\left(\sum_e C(e)t_e\right)^2\right)
\]
becomes \(\sum_e t_e=o((\sum_e t_e)^2)\), i.e.\ \(\sum_e t_e\to\infty\),
which holds because \(\sum_e t_e=(1-o(1))|Z_n|\asymp n/\log n\).  The
hypergraph is 3-bounded (each edge has size 3), \(\alpha(t)\to0\) by the
codegree-2 calculation above, and convergence in \(\alpha(t)\) is uniform in
the hypergraph by Kahn's uniformity convention.  Thus Theorem~1.5 yields
\[
  |M_n|=\sum_{A\in M_n}1\sim\sum_e t_e=(1-o(1))|Z_n|.
  \tag{5}
\]

\section{Cleanup and completion}

For each matched edge \((x,y,P)\), choose
\[
  a_P\equiv x\pmod P.
\]
Since \(|y-x|=2P\), this also hits \(y\).  The residue is nonzero: if
\(P>n/5\) divided a residual target, then the target would be one of
\(P,2P,3P,4P\), and these are nonresidual by the robust side-debt lemma.

Let
\[
  N=|A_S(n)|=(1+o(1))\frac n{\log n}.
\]
By (5) and the definition of \(Z_n\),
\[
  |M_n|=\left((\beta-\tfrac15)\delta_S+o(1)\right)N.
\]
The robust reservoir has size
\[
  |\R(n)|=\left(\frac45\delta_S+o(1)\right)N.
\]
Let \(E_S(n)=o(N)\) be the exceptional residual-token count outside the main
one-token set \(A_S(n)\).  After the pair stage, the unused robust-prime
reservoir has size
\[
  U=|\R(n)|-|M_n|,
\]
while the residual-token multiset still needing singleton cleanup has size
\[
  T_{\rm rem}=|A_S(n)|-2|M_n|+E_S(n).
\]
Thus the exact injection condition is \(U\ge T_{\rm rem}\), equivalently
\[
  |\R(n)|+|M_n|\ge|A_S(n)|+E_S(n).
\]
The left-hand side is
\(\bigl(\tfrac45\delta_S+(\beta-\tfrac15)\delta_S+o(1)\bigr)N=(1+\Delta+o(1))N\),
while the right-hand side is \(N+E_S(n)=(1+o(1))N\).  Since \(\Delta>0\) is
fixed and \(E_S(n)=o(N)\), the inequality holds for all sufficiently large
\(n\).  This is exactly~(1) translated into the cleanup ledger.

Assign each remaining residual token injectively to an unused robust prime
\(P\), choosing \(a_P\equiv x\pmod P\).  If an integer has two residual tokens,
use two distinct unused robust primes.  The robust side-debt lemma ensures that
switching these primes creates no new unresolved debt.  All unswitched primes
remain in their zero classes, \(2\) remains in the odd class, \(3\) remains in
the zero class, and the primes in \(S\) remain in the fixed nonzero classes.
Thus every integer \(m\le n\) receives at least two hits.  This completes the
proof. \qed

\section{References}

The analytic input is the Green--Tao finite-complexity theorem for systems of
affine-linear forms in primes, made unconditional using the Green--Tao
nilsequence theorem and the Green--Tao--Ziegler inverse theorem for Gowers
norms.  The original linear-equations-in-primes framework is

\begin{itemize}[leftmargin=*]
\item B. Green and T. Tao, \emph{Linear equations in primes}, Annals of
Mathematics 171 (2010), 1753--1850.
\item B. Green and T. Tao, \emph{The M\"obius function is strongly orthogonal
to nilsequences}, Annals of Mathematics 175 (2012), 541--566.
\item B. Green, T. Tao, and T. Ziegler, \emph{An inverse theorem for the
Gowers \(U^{s+1}[N]\)-norm}, Annals of Mathematics 176 (2012), 1231--1372.
\end{itemize}

The matching input is

\begin{itemize}[leftmargin=*]
\item J. Kahn, \emph{A linear programming perspective on the
Frankl--R\"odl--Pippenger theorem}, Random Structures and Algorithms 8
(1996), 149--157.
\end{itemize}

The printed Theorem~1.5 of Kahn states the pair co-load parameter
\(\alpha(t)=\max\{\sum t(A):x,y\in A\in\mathcal H\}\) over \(x\ne y\) and the
fractional rounding conclusion used in Section~7, with convergence uniform in
the hypergraph (the ``uniformity convention'' of Kahn--Pippenger--Spencer).
The 3-bounded \(C(e)\equiv 1\) application above falls inside this hypothesis.

\appendix
\section{Normalized GTZ moment input on the fixed finite core}

Freeze the global data before taking \(n\to\infty\): the finite set
\(S\subset\{7,11,13,\ldots\}\) with \(W=\prod_{s\in S}s\) and \(W_0=2W\); the
robust set \(\B\subset(\mathbb Z/W_0\mathbb Z)^\times\); the interval
\(I_\beta=(1/5,\beta]\) with \(1/5<\beta<1/2\); the finite type set
\(\mathcal T\); and bounded nonnegative kernels \(g_\tau\) on \(D_\tau\).

The prime-class scale is
\(N_{W_0}(n):=n/(\varphi(W_0)\log n)\).
For an edge \(e=(x,y,P)\) of type \(\tau\),
\(w_e=(\log^2 n/n)g_\tau(q/n,q'/n)\).  The vertex measures are
\(d\mu_Z=dt\), \(d\mu_X=dQ\), \(d\mu_Y=dQ'\), and the side-fiber masses are
\(\xi_{a,r}=1/(2a)\), \(\eta_{b,r'}=1/(2b)\).

\paragraph{Local constants \(\kappa_\tau\).}
The local constant is normalized so that for bounded piecewise-continuous
\(F\) supported in \(D_\tau\),
\[
  \sum_{\substack{q\equiv r,\,q'\equiv r'\,(W_0)\\ q,q',\sigma(bq'-aq)\text{ prime}\\
                   \sigma(bq'-aq)\equiv\pi\,(W_0)}}
  \frac{\log^2 n}{n}F(q/n,q'/n)
  =
  \bigl(N_{W_0}(n)+o(N_{W_0}(n))\bigr)
  \kappa_\tau\int_{D_\tau}F\,dQ\,dQ'.
  \tag{A.1}
\]
Locally obstructed types have \(\kappa_\tau=0\) and are removed.

\paragraph{Auxiliary \(W_{\rm GTZ}\)-tricked setup.}
Choose a small-prime cutoff \(w\), fixed while \(n\to\infty\), and set
\(W_{\rm GTZ}=\operatorname{lcm}(W_0,\prod_{\ell\le w}\ell)\).  For each
reduced \(\alpha\pmod{W_{\rm GTZ}}\), define
\[
  \Lambda_{W_{\rm GTZ},\alpha}(m)
  =\frac{\varphi(W_{\rm GTZ})}{W_{\rm GTZ}}\Lambda(W_{\rm GTZ}m+\alpha).
  \tag{A.2}
\]
Lift \(q=W_{\rm GTZ}m+\alpha\), \(q'=W_{\rm GTZ}m'+\alpha'\), with
\(\alpha\equiv r\), \(\alpha'\equiv r'\pmod{W_0}\), and the label residue
\(\alpha_P\equiv\sigma(b\alpha'-a\alpha)\pmod{W_{\rm GTZ}}\),
\(\alpha_P\equiv\pi\pmod{W_0}\).  Discard locally obstructed lifts.  The order
of limits is \(n\to\infty\) with \((\mathcal T,w,\eta)\) fixed, then
\(w\to\infty\), then the smoothing parameter \(\eta\to0\).

\paragraph{Smoothed kernels.}
Replace \(\mathbf 1_{D_\tau}g_\tau\) by a Lipschitz approximant
\(G_{\tau,\eta}=F_{\tau,\eta}g_{\tau,\eta}\) supported in an
\(\eta\)-neighborhood of \(D_\tau\), with
\(\|g_{\tau,\eta}-g_\tau\|_{L^1}\le\eta\) and
\(\|g_{\tau,\eta}\|_\infty\le\|g_\tau\|_\infty\).  Polygonal boundaries and a
finite type set make every first/second moment shift by \(o_\eta(|Z_n|)\)
uniformly.  Sharp kernels are recovered as \(\eta\to0\).

\section{Exact finite-complexity systems}

Each system below is restricted to finitely many compatible residue lifts
modulo \(W_{\rm GTZ}\) and counted using the normalized weights of (A.2).

\paragraph{B.1.\ Edge-total system.}
For type \(\tau=(a,b,\sigma,r,r',\pi)\), variables \((q,q')\), the three
prime forms
\[
  q,\qquad q',\qquad P=\sigma(bq'-aq)
  \tag{B.1}
\]
have coefficient vectors \((1,0)\), \((0,1)\), \((-\sigma a,\sigma b)\),
which are pairwise non-proportional.  Finite complexity holds.

\paragraph{B.2.\ Label second-moment system.}
For two types \(\tau_i=(a_i,b_i,\sigma_i,r_i,r_i',\pi)\) sharing a common
\(\pi\), choose integers \(u_i,v_i\) with \(b_iv_i-a_iu_i=1\) and parametrize
the common-label equation \(P=\sigma_i(b_iq_i'-a_iq_i)\) by \((P,t_1,t_2)\):
\[
  q_i=u_i\sigma_iP+b_it_i,\qquad
  q_i'=v_i\sigma_iP+a_it_i.
\]
The five prime forms
\[
  P,\quad u_1\sigma_1P+b_1t_1,\quad v_1\sigma_1P+a_1t_1,\quad
  u_2\sigma_2P+b_2t_2,\quad v_2\sigma_2P+a_2t_2
  \tag{B.7}
\]
have coefficient vectors
\(
  (1,0,0)
\), \(
  (\sigma_1u_1,b_1,0)
\), \(
  (\sigma_1v_1,a_1,0)
\), \(
  (\sigma_2u_2,0,b_2)
\), \(
  (\sigma_2v_2,0,a_2)
\)
in \((P,t_1,t_2)\).  Pairwise non-proportionality holds: forms involving
distinct \(t_i\)'s are evidently independent, and within a fixed \(i\),
\((\sigma_iu_i,b_i)\propto(\sigma_iv_i,a_i)\) would force
\(a_iu_i=b_iv_i\), contradicting \(b_iv_i-a_iu_i=1\).

\paragraph{B.3.\ \(X\)-side second-moment system.}
Two edges share an \(X\)-vertex only when \(a_1q_1=a_2q_2\); for large \(n\)
this forces \(a_1=a_2=a\), \(q_1=q_2=q\), \(r_1=r_2=r\).  In variables
\((q,q_1',q_2')\), the five prime forms are
\[
  q,\quad q_1',\quad q_2',\quad
  \sigma_1(b_1q_1'-aq),\quad
  \sigma_2(b_2q_2'-aq),
  \tag{B.13}
\]
with coefficient vectors
\(
  (1,0,0),(0,1,0),(0,0,1),(-\sigma_1a,\sigma_1b_1,0),(-\sigma_2a,0,\sigma_2b_2)
\),
pairwise non-proportional in the ambient 3-variable system.  Collision
slices such as \(q_1'=q_2'\) are lower-dimensional and removed in
Appendix~C.

\paragraph{B.4.\ \(Y\)-side second-moment system.}
Symmetric to B.3.  Shared \(Y\)-vertices force \(b_1=b_2=b\),
\(r_1'=r_2'=r'\), \(q_1'=q_2'=q'\).  In variables \((q',q_1,q_2)\), the five
prime forms
\[
  q',\quad q_1,\quad q_2,\quad
  \sigma_1(bq'-a_1q_1),\quad
  \sigma_2(bq'-a_2q_2)
  \tag{B.18}
\]
have coefficient vectors
\(
  (1,0,0),(0,1,0),(0,0,1),(\sigma_1b,-\sigma_1a_1,0),(\sigma_2b,0,-\sigma_2a_2)
\),
pairwise non-proportional away from the lower-dimensional collision slices.

\section{Diagonal and exceptional-set deletion}

All deletions are uniform over the finite type set; since
\(|Z_n|\asymp n/\log n\), it suffices that each deleted weighted contribution
is \(o(n/\log n)\).

\paragraph{C.1.\ Identical-edge diagonal.}
Per type, \(O(n^2/(\log n)^3)\) edges, each weighted by
\(O(\log^4 n/n^2)\), giving a total identical-edge diagonal of
\(O(\log n)=o(|Z_n|)\).

\paragraph{C.2.\ Other form-collision slices.}
Slices like \(q_1=q_2\) or \(q_1'=q_2'\) in second-moment systems have at most
\(O(n^2)\) integer points before primality, so weighted total
\(O(\log^4 n)=o(|Z_n|)\).  After deletion the remaining systems have pairwise
non-proportional affine forms.

\paragraph{C.3.\ Small-prime values.}
For fixed \(w\), removing tuples in which a prime form is \(\le w\) reduces
the configuration count by \(O_w(n^2)\) at the cost of weighted contribution
\(O_w(\log^4 n)=o(|Z_n|)\).  This justifies passing to reduced residue classes
modulo \(W_{\rm GTZ}\).

\paragraph{C.4.\ Boundary layers.}
For fixed \(\eta\), the \(\eta\)-neighborhood of \(\partial D_\tau\) has area
\(O_\tau(\eta)\); first moments lose \(O_\tau(\eta)N_{W_0}(n)\), second
moments lose \(O_{\tau_1,\tau_2}(\eta)|Z_n|\).  Letting \(\eta\to0\) after
\(n\to\infty\) recovers the sharp kernels.

\section{Local admissibility}

\paragraph{D.1.\ Fixed small primes in \(W_0=2W\).}
Retained residue classes enforce \(q,q',P\in(\mathbb Z/W_0\mathbb Z)^\times\),
so no \(\ell\mid W_0\) divides any prime form on a retained class.  At
\(\ell=2\), the constraint \(bq'-aq\equiv-q\pmod 2\) is odd because \(a\) is
odd, \(b\) is even, and \(q\) is odd in retained classes.  At \(\ell\in S\),
the unit \(\pi\) forbids \(\ell\mid a\) and \(\ell\mid b\) simultaneously
(else \(\sigma(br'-ar)\equiv 0\pmod\ell\)), so retained types automatically
satisfy \(\gcd(a,b)=1\); the remaining cases give a unit residue for
\(P\pmod\ell\).

\paragraph{D.2.\ Small primes introduced by \(W_{\rm GTZ}\).}
For \(\ell\le w\), \(\ell\nmid W_0\), retain a residue lift only when every
prime form in the relevant system is a unit modulo \(\ell\).  This is a
finite check for each \((w,\mathcal T)\).

\paragraph{D.3.\ Large primes.}
For \(\ell\nmid W_{\rm GTZ}\), after the deletions of Appendix~C no two
surviving forms are rationally affinely dependent.  For all sufficiently large
\(\ell\), the union of at most five hyperplanes \(L_i=0\) cannot cover
\(\mathbb F_\ell^d\), giving a point where all forms are nonzero.

\paragraph{D.4.\ Local-factor identities.}
In the normalized \(W_{\rm GTZ}\)-tricked formulation, each one-form
density has been normalized to mean \(1\) in its lifted residue class, so the
joint two-edge densities factor as
\(
  \kappa^Z_{\tau_1,\tau_2}=\kappa_{\tau_1}\kappa_{\tau_2}
\),
\(
  \kappa^X_{\tau_1,\tau_2}=\kappa_{\tau_1}\kappa_{\tau_2}/\xi_{a,r}
\),
\(
  \kappa^Y_{\tau_1,\tau_2}=\kappa_{\tau_1}\kappa_{\tau_2}/\eta_{b,r'}
\)
automatically.  These are the Euler-factor conditional-density identities for
the joint system: density of joint = product of edge densities divided by
density of shared vertex.

\section{Disintegration of main terms; absence of pointwise prime-pair input}

\paragraph{E.1.\ What GTZ is asked to prove.}
For each fixed \((w,\eta)\), each type or type pair, and each residue lift,
the normalized \(W_{\rm GTZ}\)-tricked theorem yields
\[
  \sum_{\mathbf m\in K_n\cap\mathbb Z^d}
    \prod_{j=1}^k \Lambda_{W_{\rm GTZ},\alpha_j}(L_j(\mathbf m))
    \Phi(\mathbf m/n)
  =\operatorname{Vol}_\Phi(K_n)+o_{w,\eta}(n^d),
\]
with \(\Phi\) the appropriate bounded Lipschitz kernel product.

\paragraph{E.2.\ Label disintegration.}
The off-diagonal label second-moment main term for \(\tau_1,\tau_2\) is
\[
  N_{W_0}(n)\int_{I_\beta}
    L_{Z,\tau_1}^{\rm lim}(\pi,t)L_{Z,\tau_2}^{\rm lim}(\pi,t)\,dt
  +o(|Z_n|).
\]
Summing over types and using \(L_Z^{\rm lim}=1\) gives
\(\sum_P L_Z^2=|Z_n|+o(|Z_n|)\); combined with the edge-total system,
\(\sum_P(L_Z-1)^2=o(|Z_n|)\).  This is an averaged 3-variable count in
\((P,t_1,t_2)\), \emph{not} a pointwise estimate for primes representing a
fixed \(P\).

\paragraph{E.3--E.4.\ Side disintegrations.}
Identical structure to E.2 but with the shared vertex playing the role of
\(P\); the result is the \(X\)- and \(Y\)-side \(L^2\) bounds claimed in the
weighted moment proposition.

\paragraph{E.5.\ Why no hidden pointwise input.}
The proof never asks for an asymptotic for primes \(q,q'\) with
\(bq'-aq=P\) for fixed \(P\), nor for pointwise side degrees, nor for any
Hardy--Littlewood/Bateman--Horn singular-series estimate varying a single
fixed label.  Every second moment averages over the shared vertex; these are
the multidimensional finite-complexity systems standardly covered by GTZ.

\section{Final proposition for the moment block}

\begin{proposition}[Fixed-core normalized GTZ weighted moments]\label{prop:F1}
With \(S,W,W_0,\B,\beta,\mathcal T\), kernels \(g_\tau\), and constants
\(\kappa_\tau\) fixed as above, assume the normalized
\(W_{\rm GTZ}\)-tricked Green--Tao--Ziegler finite-complexity theorem for
the systems in Appendix~B.  Then, as \(n\to\infty\),
\[
  \sum_{P\in Z_n}L_Z(P)=|Z_n|+o(|Z_n|),
  \quad
  \sum_{P\in Z_n}(L_Z(P)-1)^2=o(|Z_n|),
\]
\[
  \sum_{x\in X_n}(L_X(x)-L_X^{\rm lim}(x))^2=o(|Z_n|),
  \quad
  \sum_{y\in Y_n}(L_Y(y)-L_Y^{\rm lim}(y))^2=o(|Z_n|).
\]
The error terms are uniform over the fixed finite type set.  Auxiliary
parameters are removed in the order \(n\to\infty\), then \(w\to\infty\), then
\(\eta\to0\).
\end{proposition}

This is the input to fractional rounding in Section~7, together with the
limiting slack \(L_X^{\rm lim},L_Y^{\rm lim}\le1-2\gamma\) and the atom bound
\(\max_e w_e\ll(\log^2 n)/n=o(1)\).

\begin{thebibliography}{9}
\bibitem{Kahn1996}
J.\ Kahn, \emph{A linear programming perspective on the
Frankl--R\"odl--Pippenger theorem}, Random Structures and Algorithms 8
(1996), 149--157.
\end{thebibliography}

\end{document}
